\subsection{Five-Cycle Structure and Odd-Degree Reductions}
\label{app:multi-structure}

Use the notation of \textbf{Section~\ref{sec:exact}}. The host is
induced by the endpoints of the marked copies, and $D$ remains the
degree bound after this restriction. Write $\mu=\frac{5}{4}D^2-m$.
The capacity and redistribution estimates below apply to both parities
of $D$. When we apply them to a counterexample of odd degree, we use
$\mu\le\frac{D}{2}-\frac{5}{4}$.

\subsubsection{Core capacities and discharging}

Let $c_0,\ldots,c_4$ be the vertices of an induced cycle in the simple
support whose five edges are red. Call these vertices the anchors.
Indices in this appendix are taken modulo five. The class $U_i$ consists
of the vertices whose anchor neighbors are exactly $c_{i-1}$ and
$c_{i+1}$, with both pairs red. Write $B=\bigcup_i U_i$ for the
original core. A vertex outside $B$ is called an exterior vertex
until the parts are enlarged. Let $L$ count the blue edges between $U_i$ and
$U_{i+1}$ over all $i$. Every vertex of $U_i$ is adjacent in the
simple support to every vertex of $U_{i+1}$, as proved in
\textbf{Section~\ref{sec:exact}}. Put
$a_i=|U_i|$, $s_i=D-a_{i-1}-a_{i+1}$, $R=\sum_i s_i$ and
$a_*=\frac{D-R}{2}$. The quantity $s_i$ is an upper bound on the
number of host incidences at a vertex of $U_i$ after reserving one edge
to each vertex of $U_{i-1}\cup U_{i+1}$. Summing the defining equations
gives $|B|=\frac{5D-R}{2}$. Since
$s_{i-2}+s_{i+2}=2D-|B|+a_i$, we obtain
\begin{equation}\label{eq:multi-core-identities}
a_i=a_*+s_{i-2}+s_{i+2},\qquad
A=\sum_i a_i s_i=\frac{RD}{2}-\frac{R^2}{2}
                    +2\sum_i s_i s_{i+2}\le\frac{RD}{2}.
\end{equation}
For the inequality, take the five-cycle on the indices with edges
$\{i,i+2\}$. If two nonadjacent indices have positive weights, move
their combined weight to the index with the larger sum of neighboring
weights. The product sum weakly increases and the support shrinks.
After repeating the move, the support is a clique of size at most two,
so the product sum is at most $\frac{R^2}{4}$. The identity
$R=5D-2|B|$ gives $R\equiv D\pmod 2$.

The signature of a vertex $x$ is the set $\{i:xc_i\in E(G_{\rm s})\}$,
of size $k_x$. Summing the capacity $s_i$ at the five anchors gives
$K=\sum_{x\notin B}k_x\le R$. For $x\notin B$, let $h_x$ be its marked
degree, let $Q_x=|N_{G_{\rm s}}(x)\cap B|$, and let $b_x$ count its
blue pairs to $B$. For any marked edge $xy$, the union of the
signatures of $x$ and $y$ covers every edge of the cycle. In fact,
if $xy$ is disjoint from a cycle edge, the strong clique condition
gives an adjacency to one of the ends of that cycle edge. If $xy$
shares an endpoint with a cycle edge, its other cycle endpoint is
a neighbor of the shared vertex. Since a vertex cover of $C_5$ has
at least three vertices, a vertex with $k_x=0$ can have a marked
edge only to a vertex $y$ with $k_y\ge3$. In particular such an $x$
has no marked edge into $B$, since $k_y=2$ for $y\in B$.
A vertex outside $B$ of signature size two has at least one blue anchor pair. For
consecutive positions, a marked anchor pair would fail the strong
clique condition with the cycle edge opposite those positions. For
nonconsecutive positions, two marked anchor pairs would put the vertex
in $B$.

Every blue pair between consecutive classes uses two unmarked incidences
in $B$. The $b_x$ blue pairs of $x$ into $B$ use one unmarked incidence
in $B$ each, disjoint from the first set of incidences. The two
counts are illustrated in \textbf{Figure~\ref{fig:multigraph-cut}}.
Summing the marked incidences inside and outside $B$ gives
\begin{equation}\label{eq:multi-core-cut}
 2m\le D|B|-2L+\sum_{x\notin B}(h_x-b_x).
\end{equation}

\begin{figure}[htbp]
\centering
\begin{tikzpicture}[
  marked edge/.append style={draw=red!70!black},
  blue pair/.style={draw=blue!70!black,densely dashed,line width=0.95pt},
  font=\small
]
  \filldraw[fill=black!4,draw=black!40,rounded corners=4pt]
    (-0.65,-0.55) rectangle (4.15,1.85);
  \node[anchor=north west] at (-0.65,1.85) {$B$};
  \node at (0.15,1.51) {$U_i$};
  \node at (3.3,1.51) {$U_{i+1}$};
  \node[graph vertex] (a) at (0.15,0.95) {};
  \node[graph vertex] (b) at (0.15,0) {};
  \node[graph vertex] (c) at (3.3,0.95) {};
  \node[graph vertex] (d) at (3.3,0) {};
  \node[graph vertex,label=above:$x$] (x) at (5.95,0) {};
  \draw[marked edge] (a) to[bend left=9] (c);
  \draw[marked edge] (a) to[bend right=9] (c);
  \draw[marked edge] (a)--(d);
  \draw[marked edge] (b)--(c);
  \draw[blue pair] (b)--node[below] {two incidences} (d);
  \draw[blue pair] (d)--node[above] {one incidence} (x);
\end{tikzpicture}
\caption{A blue edge between consecutive parts contributes twice to
the degree sum on $B$. A blue edge from $x\notin B$ to $B$ contributes
once. Both copies of a marked pair contribute to marked degree.}
\label{fig:multigraph-cut}
\end{figure}


We first bound the total outside contribution after charging $b_x$ against
the core capacity. Suppose $k_x=1$ and $x$ has a marked edge to a vertex of $B$.
We show that $b_x\ge a_*+1$. Let $j$ be the sole index in the signature of $x$, and write $xy$ for
its marked edge with $y\in U_i$. Since $\{j,i-1,i+1\}$ covers the
cycle, $i=j\pm2$. Of the two cycle neighbors of $c_j$, let $c_k$
be the one nonadjacent to $y$. For each $u\in U_j$ the edge $c_ku$
is marked. Comparing it with $xy$ gives $xu$ or $yu$ in the support.
The edge $xy$ uses one of the $s_i$ extra incidences at $y$, so
$y$ has at most $s_i-1$ neighbors in $U_j$. At least
$a_j-s_i+1\ge a_*+1$ vertices of $U_j$ are therefore adjacent to
$x$. Each of these edges is blue, since $\{j,j-1,j+1\}$ does not
cover the cycle. If $xy$ is marked with $k_x=1$ and $k_y=2$, the union of their
signatures is a vertex cover of the cycle and hence has size three.
Whenever the union on an outside marked edge has size three, its two
omitted positions are nonconsecutive. Every red edge from either
omitted anchor into $B$ forces its other end to be adjacent to at least
one end of the outside marked edge. The core neighborhoods of those
ends therefore cover three classes. The other two classes have total
size at most $D$, so the three classes contain at least $|B|-D=D+a_*$
vertices. Each end uses one degree unit on the outside marked edge,
and so each sees at least $a_*+1$ vertices of $B$. If the end of
signature size one has no marked edge into $B$, these pairs are blue.

We redistribute the quantities $h_x-b_x$ over vertices outside $B$.
A vertex of signature size one with a marked core neighbor or a
marked neighbor of signature size two has $b_x\ge a_*+1$. Since
$h_x+b_x\le D$, it retains a contribution at most $R-2$. Any other
vertex of signature size one sends $\frac{1}{2}$ per marked copy to
its neighbors, all of which have signature size at least three. Its
remaining contribution is at most $\frac{D}{2}$. A vertex $x$ of
signature size zero with a marked neighbor of signature size three
has $b_x=Q_x\ge a_*$ by the comparison with the two omitted anchors.
It distributes its positive contribution uniformly over its marked
edges. Since $h_x\le D-Q_x$, the amount per copy is at most
$\frac{2R}{D+R}$. The other vertices of signature size zero use the
same rule, with at most one per copy and with recipients of
signature size at least four. We discard any negative contribution
left at a vertex of signature size zero.

A recipient of signature size three with a marked neighbor of
signature size zero also sees at least $a_*$ vertices of $B$, so it
has at most $\frac{D+R}{2}$ marked copies outside $B$. Multiplying
this bound by the larger of the two sending rates shows that it
receives at most $\max\{\frac{D+R}{4},R\}$. If it has no marked
neighbor of signature size zero, it receives at most $\frac{D}{2}$.
A recipient with at least four positions receives at most $D$.
Adding each recipient's own contribution, at most $D$, shows that
the total contribution at every outside vertex is at most
$k_x\max\{\frac{D}{2},R\}$ for $R\le D$. By \eqref{eq:multi-core-cut},
\begin{equation}\label{eq:multi-rough-discharge}
2m\le D|B|-2L+\max\left\{\frac{D}{2},R\right\}K.
\end{equation}

In the graph on $B$ consisting of the marked pairs, the minimum size
of a vertex cover is at least $D+q$, where $q=a_*-\frac{L}{a_*}$
when $a_*>0$. In fact, any independent set in this graph with $x_i$
vertices in $U_i$ satisfies $\sum_i x_i x_{i+1}\le L$, and
\[
 \sum_i x_i x_{i+1}\ge a_*\left(\sum_i x_i-D\right).
\]
The difference between the two sides is affine in each $x_i$ on
$[0,a_i]$, so it suffices to check the inequality for each $x_i$ in
$\{0,a_i\}$. At such a point, delete one occupied class from each
consecutive occupied pair until the occupied classes are pairwise
nonconsecutive. An edge with both classes occupied accounts for a
product $a_i a_{i+1}$, and either deletion removes at most
$\frac{a_i a_{i+1}}{a_*}$ vertices. The remaining classes contain at
most $D$ vertices because any two nonconsecutive classes have a
common neighboring class. Hence any independent set has size at most
$D+\frac{L}{a_*}$. Subtracting from $|B|=2D+a_*$ gives the bound on
the vertex cover number. The argument applies also with enlarged
classes below. The core neighborhoods of the two ends of an outside
marked edge cover all marked pairs in $B$, so $Q_x+Q_y\ge D+q$ for
every such edge $xy$.

For a sharper estimate, let $q\ge0$ be any lower bound on the vertex
cover number minus $D$. Vertices of signature size zero, and vertices
of signature size one with no marked neighbor in $B$ or of signature
size two, now send only their positive contributions. They distribute
these contributions uniformly over their marked copies, all of which
meet vertices of signature size at least three. At a sender $x$ with
positive contribution, $b_x=Q_x$, so $Q_x<h_x\le D-Q_x$.
Each of its marked neighbors has at least $D+q-Q_x$ core neighbors
and hence at most $Q_x-q$ outside marked copies. In particular,
$q\le Q_x<\frac{D}{2}$. If $q\ge\frac{D}{2}$, there is no positive
transfer. Set $u=(\frac{D}{2}-q)_+$. A recipient receives at most
\[
 C(u)=\frac{2u^2}{\frac{D}{2}+u}.
\]
Indeed, the amount sent per copy is
$1-\frac{Q_x}{h_x}\le\frac{D-2Q_x}{D-Q_x}\le\frac{D-2q}{D-q}$,
and the recipient has at most $u$ outside marked copies. If $a_*>0$
and $q\ge a_*-\frac{L}{a_*}$, then $u\le\frac{R}{2}+\frac{L}{a_*}$.
Since $0\le C'(u)\le2$ for $u\ge0$, it follows that
$C(u)\le\frac{R^2}{D+R}+\frac{2L}{a_*}$.
More generally, if there are $c$ possible recipients, positive
transfer requires $Q_x<h_x\le c(Q_x-q)$. This is impossible for $c=1$,
and for $c=2$ it implies $Q_x>2q$. For two recipients the amount
sent per copy is at most $\frac{D-4q}{D-2q}$ whenever transfer
is possible. Multiplying by $u$ shows that each recipient receives
at most $\max\{0,\frac{D-4q}{2}\}$.

For each $x$ outside $B$, discard any negative contribution left
at a sender and denote by $\ell_x$ the difference between
$\frac{k_xD}{2}$ and the resulting contribution of $x$.
A vertex of signature size one has $\ell_x\ge t+2$ for
$t=\frac{D}{2}-R$ and $R\ge3$. Indeed, it either retains at most
$R-2$ or sends its whole positive contribution and retains zero.
For signature size two, $\ell_x=D-h_x+b_x\ge2b_x\ge2$.
A recipient of signature size $k$ has
$\ell_x\ge(\frac{k}{2}-1)D-C(u)$, or the bound obtained using the
smaller transfer estimate when $c\le2$. Substituting
$\sum_x(h_x-b_x)\le\frac{DK}{2}-\sum_x\ell_x$ into
\eqref{eq:multi-core-cut} and using $|B|=\frac{5D-R}{2}$ gives
\begin{equation}\label{eq:multi-loss-budget}
 2L+\frac{D}{2}(R-K)+\sum_x\ell_x\le2\mu.
\end{equation}

\subsubsection{A joint bound on the core and its exterior}
\label{app:multi-joint}

Write $r=\frac{R}{D}$, $j=\frac{\sum_i j_{c_i}}{D}$,
$g=\frac{m}{D^2}-\frac{5}{4}$, and $\lambda=\frac{L}{D^2}$.
The outside anchor pairs and the extra marked copies at the anchors
use disjoint units of the capacity $R$, so $\frac{K}{D}\le r-j$.
If $\frac{1}{2}\le r\le1$, then
\begin{equation}\label{eq:multi-joint}
 2g+2\lambda\le
 \frac{1+r}{2+r}\,r\left(r-\frac{1}{2}\right)
 -\frac{1+r^2}{2+r}\,j .
\end{equation}

To prove this, divide all vertex counts and degrees by $D$, and all
edge counts by $D^2$. Write $\eta=\frac{1-r}{2}$, $v=\frac{1+r}{2}$
and $\rho=\frac{2r}{1+r}$. An outside marked edge whose signatures
have union of size three has two ends each with normalized core
degree at least $\eta$. Each end has outside marked degree at most
$v$.

Let $x$ count the vertices of signature size one with a marked
neighbor in the core or a marked neighbor of signature size two.
Their total contribution in \eqref{eq:multi-core-cut}, in the
normalized units, is denoted by $C_1$. Each of these vertices has
at least $\eta$ blue pairs into the core, and $C_1\le rx$. Let
$p_2,p_3$ count the marked edges from these vertices to vertices of
signature size two and of signature size at least three, respectively.
Let $n_2$ count their neighbors of signature size two via a marked
edge. Also let $z_3,z_0$ count the vertices of signature size three
and of signature size zero, respectively, that meet a marked edge
joining these two groups. Denote the count of such edges by $p_{03}$.
The three groups counted by $n_2,z_3,z_0$ are disjoint. Each of their
vertices uses at least $\eta$ of the core capacity by the comparison
with the two omitted anchors, and $p_2\le v n_2$ and
$p_{03}\le v\min\{z_0,z_3\}$.

By \eqref{eq:multi-core-identities}, the total normalized capacity
of the core after reserving its consecutive pairs is at most
$\frac{r}{2}$. Deduct the core capacity used by the groups counted
by $x,n_2,z_3,z_0$. To bound $C_1$, add the exterior marked
counts $p_2+p_3$ and deduct the blue core pairs at the vertices
counted by $x$. This gives
$C_1\le\frac{r}{2}-2\eta x-\eta(n_2+z_3+z_0)+p_2+p_3$.
With $\theta=\frac{2r-1}{2+r}$ and $c=r-\theta=\frac{1+r^2}{2+r}$,
take weights $1-\theta,\theta$ on the two bounds for $C_1$ to get
\begin{equation}\label{eq:multi-joint-first}
 C_1\le cx+\frac{\theta r}{2}
     +\theta(p_2+p_3)-\theta\eta(n_2+z_3+z_0).
\end{equation}

For the other outside vertices, use the redistribution leading to
\eqref{eq:multi-rough-discharge}. A vertex of signature size two
retains at most one. Assign $\theta p_2-\theta\eta n_2$ in
\eqref{eq:multi-joint-first} to the group of such vertices. Since
$p_2\le v n_2$, its additional contribution is at most $\theta r$
per vertex, or at most $c=\frac{1}{2}+\frac{\theta r}{2}$ per signature
position after adding the original contribution.

For a vertex of signature size at least four, let $d_1$ be its
marked degree to the group counted by $x$. None of those edges
transfers any contribution to it, so its own and received
contributions sum to at most $2-d_1$. The corresponding bound for a
vertex of signature size three with no marked neighbor of signature
size zero is $\frac{3}{2}-\frac{d_1}{2}$. Adding the corresponding
$\theta d_1$ from $\theta p_3$ gives at most $c$ per signature
position in both cases, since $c\ge\frac{1}{2}$ and $\theta\le\frac{1}{3}$.

For the group counted by $z_3$, write $d_{1,3}$ for its marked
degree to the group counted by $x$. The vertices counted by $z_3$
have outside marked degree at most $v$ and receive at most
$\frac{1}{2}$ per marked copy, except for edges to the vertices
counted by $z_0$, which can transfer up to $\rho$ each. The total
after adding the corresponding terms of \eqref{eq:multi-joint-first}
is at most
\[
 \left(1+\frac{v}{2}\right)z_3
 -\left(\frac{1}{2}-\theta\right)d_{1,3}
 +\left(\rho-\frac{1}{2}\right)p_{03}-\theta\eta(z_3+z_0).
\]
Use $z_0\ge\frac{p_{03}}{v}$, then $p_{03}\le v z_3$, to
bound this by $(1+\rho v-2\theta\eta)z_3=3c z_3$. The second
substitution is valid because
$\rho-\frac{1}{2}-\frac{\theta\eta}{v}\ge0$ for $\frac{1}{2}\le r\le1$.
This accounts for all terms of \eqref{eq:multi-joint-first}.
The vertices of signature size one outside the group counted by
$x$ retain at most $\frac{1}{2}$ each. Thus the right side of
\eqref{eq:multi-core-cut}, divided by $D^2$, is at most
$\frac{5-r}{2}-2\lambda+c(r-j)+\frac{\theta r}{2}$.
Subtracting $\frac{5}{2}$ proves \eqref{eq:multi-joint}.

\subsubsection{Enlarging the Five Parts}

Every vertex $x$ whose signature consists of two nonconsecutive indices $i,j$,
with $b_x<a_*$, and with complete host adjacency to both $U_i$ and
$U_j$, is put into their common neighboring class. We show that
every pair joining consecutive enlarged classes is still an edge
of the simple support. For two newly
added vertices $x,y$ in consecutive classes $p,p+1$, choose a marked
neighbor of $x$ in $U_{p-1}$ avoiding the at most $s_{p+1}$ neighbors
of $y$ there. Then choose a marked neighbor of $y$ in $U_{p+2}$ avoiding
the neighbors of the first choice and of $x$. The latter two vertices
have at most $s_{p-1}+s_p$ neighbors in that class. Both choices
are possible by \eqref{eq:multi-core-identities} and $b_x,b_y<a_*$.
The two chosen marked edges force $xy$. The old--new pairs are
complete by the inclusion rule. A prime denotes enlarged quantities.
In particular each inclusion reduces $R$ and $K$ by two. The old
marked anchor stars can still be used.

Here are the comparisons for a remaining vertex of signature size two. If its
positions $\{i,j\}$ are nonconsecutive, $b_x<a_*$, and
it remains outside, it omits at least one point of both $U_i$ and $U_j$.
Indeed an omitted point in $U_i$, through its marked spoke, bounds the
number of distinct marked neighbors of $x$ in $U_j$ by $s_i$. This
contradicts complete adjacency to $U_j$ since $a_j-b_x>s_i$.
The same comparison with an omitted point works for neighbors in
$U'_j$, using the new bound $s_i'$. Its marked core neighbors can only
lie in $U_i'\cup U_j'$. To see this, take positions $\{4,1\}$, so that
comparison with the cycle excludes $U_0'$. A mark to $U_2'$ would force
at least $a_4-s_2'$ host neighbors of $x$ in $U_4$ by the star from
$c_0$. At most $s_1'$ of these are marked, contrary to $b_x<a_*$
by \eqref{eq:multi-core-identities}. A mark to $U_3'$ is excluded by
reflection. Both anchor pairs of $x$ are blue. If one were marked,
the marked star into the other designated class, based at an anchor
missing the first, would force complete adjacency to that other
class. Each of the two marked-neighbor counts into the designated
classes is at most the slack of the opposite class. Thus
\begin{equation}\label{eq:multi-low-blue-outside}
 Q_x\le s_i'+s_j'+b_x.
\end{equation}
If $xy$ is marked with $y\in U_i'$, every omitted old $U_j$ vertex must
see $y$, again by its marked spoke. At most $s_i'-1$ such vertices
can be omitted since $xy$ already consumes one unit at $y$. It follows
that $b_x\ge a_j-2s_i'+1\ge a_*-s_i'+1$. The same omitted vertices
also give
\begin{equation}\label{eq:multi-rectangle}
 I+b_{x,j}\ge a_j-s_i'\ge a_*,
\end{equation}
where $I$ counts host edges inside the enlarged core other than its
reserved base edges and
$b_{x,j}$ counts the blue neighbors of $x$ in old $U_j$.

A vertex with $b_x<a_*$ and signature $\{0,1\}$ can have marked edges
into the enlarged core only to $U_3'$. The cycle excludes marked edges to $U_0'$ and
$U_1'$. A mark to $U_2'$ forces at least $a_0-s_2'$ blue neighbors
in $U_0$ by the star from $c_4$, and a mark to $U_4'$ is excluded
by the reflected comparison. If there is a marked neighbor in $U_3'$,
the two ends of that mark must cover the marked $U_0$--$U_1$ pairs.
If the union of their neighborhoods omits $t_0$ vertices of $U_0$
and $t_1$ vertices of $U_1$, then all $t_0t_1$ pairs between these
omitted sets are blue. Thus $t_0t_1\le L$ and one of $t_0,t_1$ is at
most $\sqrt L$.
The neighbor in $U_3'$ can pay for at most $s_3'-1$ wrong-class
adjacencies in either part after accounting for the selected mark.
Thus $b_x\ge a_*+1-\sqrt L$. Finally, if a remaining outside vertex
has no enlarged core mark, a marked outside neighbor has $Q\le D-1$.
The old vertex-cover bound therefore gives $b_x=Q_x\ge q+1$.

Reserve one base copy for each pair in consecutive enlarged
classes, choosing a marked copy whenever possible. Let $L'$
count the unmarked base copies. Let $I$ count the other host
edges inside the enlarged core, $P$ count the marked ones among
them, $L_{\rm cut}$ count the unmarked copies across the cut of
the enlarged core, and $f$ count the marked copies with both ends
outside. There are $\frac{D|B'|-A'}{2}-L'$ marked base copies.
The number of marked copies across the cut is at most
$A'-2I-L_{\rm cut}$, since the edges counted by $I$ each use two
units of the remaining core capacity $A'$. Adding $P$ and $f$ and
using \eqref{eq:multi-core-identities} gives
\begin{equation}\label{eq:multi-extra-capacity}
 \mu\ge \frac{R'D-2A'}{4}+L'+(2I-P)+L_{\rm cut}-f
 \ge L'+I+L_{\rm cut}-f .
\end{equation}
When $R'=1$, all but one class have zero extra capacity. If all
outside signatures are singletons, outside marks are impossible.
Any non-base mark in the only positive-capacity class or from
outside to that class misses the opposite cycle edge, so only marked
base copies remain. Their total is at most $f(D)$.

\subsubsection{The half-degree threshold}
\label{app:multi-half}

We prove that for odd $D\ge9$, a counterexample has $R>\frac{D+1}{2}$ on
every anchored cycle. If $R=1$, outside signatures have size at most
one and no outside vertex can have signature zero in the marked
endpoint host. The $R'=1$ base count above applies without enlargement.
Assume $R\ge3$ and $t\ge-\frac{1}{2}$. If $t\ge\frac{1}{2}$,
\eqref{eq:multi-rough-discharge} gives $L\le \frac{D}{2}-\frac{5}{4}$ and $R-K\le1$.
Also $a_*\ge3$ and $\frac{L}{a_*}<2$. For $t=-\frac{1}{2}$ and $D\ge13$, it gives
$L\le\frac{5D-9}{8}$ and $\frac{L}{a_*}<\frac{5}{2}$, with $a_*\ge3$. In both cases
$q>0$. The remaining case $(D,R)=(9,5)$ has $a_*=2$, $L\le4$ and
$q\ge0$. Also $K\le5$, so at most one recipient is possible there.
In all cases $C(u)\le \frac{D}{2}$ and the losses in
\eqref{eq:multi-loss-budget} are nonnegative. That inequality gives
$R-K\in\{0,1\}$ and $L\le \frac{D}{2}-\frac{5}{4}$.

There is at most one vertex of signature size at least three.
Indeed, $\frac{R^2}{D+R}\le \frac{D}{6}$ for $t\ge\frac{1}{2}$ and is at most $\frac{D}{6}+\frac{1}{2}$
for $t=-\frac{1}{2}$. The loss of three recipients, with $2L$ included,
would be at least $D-\frac{3}{2}$, since $a_*\ge3$. With two recipients,
the improved bound gives total loss with $2L$ at least $D-1$.
Both bounds exceed the budget. With one recipient no positive
transfer occurs, and its size is exactly three.

Suppose first there is one vertex of signature size three, say $z$. Let $a$
count vertices of signature size one. Parity and the budget force $a$ even
and $R-K=0$. The lower bound on the total loss is
$\frac{D}{2}+R-3+a(t+1)$. For $a\ge2$ this exceeds $D-\frac{5}{2}$, so $a=0$.
We show that every vertex with $k_x=2$ is included in the enlarged core.
For $t=-\frac{1}{2}$ the budget forces $L=0$ and each vertex of signature size two
to have one blue neighbor. A remaining vertex with nonconsecutive
signature positions and $b_x<a_*$ would have two blue anchor pairs.
The other remaining types
violate $b_x\ge a_*+1-\sqrt L$ or $b_x\ge q+1$. For $t\ge\frac{1}{2}$ a
remaining vertex with nonconsecutive signature positions, $b_x<a_*$
and a marked edge to the core satisfies
$b_x\ge a_*+1-s_i'\ge t+\frac{3}{2}$. Indeed, $R'=K'$ and the slack
at each position counts the remaining vertices that use it.
There are $\frac{R'-3}{2}$ such vertices of signature size two,
and at most one more can use that position, namely $z$.
Together with the loss of $z$ and the other vertices of signature size two
this gives a loss of at least $D+t-2>D-\frac{5}{2}$. A remaining vertex with $b_x\ge a_*$
costs at least $2a_*$ and a remaining vertex of either other type
costs at least $2a_*+2$ with $2L$ included. These also exceed the
budget after adding the other vertex losses.

The resulting enlarged core has $R'=3$, with three units at the
positions of $z$. Write $D=2d+1$. If the omitted positions are
nonconsecutive, the three classes of positive capacity have sizes
$d+1,d,d$, base count $f(D)-D$ and extra capacity $3d+1$.
If $zy$ is marked outside, then the vertex of signature size zero $y$ has
only blue core pairs. The two omitted anchor stars force
$Q_z'+Q_y'\ge3d+1-L'$, since each omitted point of those three
classes consumes a distinct blue base pair. Also $f\le D-Q_z'$.
Thus $L'+L_{\rm cut}-f\ge d$. If there is no outside mark, only
$z$ remains outside. The extra marks either all touch $z$ and
number at most $D$, or there is a marked non-base copy inside
the core. No same-class mark is possible by the opposite cycle
edge. A mark from the class of size $d+1$ to one of size $d$
blocks any mark from $z$ to those two classes. Its endpoints use
their only extra units and every possible edge joining their marked edge to
such a mark from $z$ would use one of those units. The total number
of extra marks is then at most $\frac{3d+1+d}{2}\le D$. If there are
only repeated internal marks between the two consecutive size-$d$
classes, one of those repetitions blocks marks from $z$ to the
size-$(d+1)$ class by the same comparison. The extra marks then
use at most $2d$ units and still number at most $D$.

If the two omitted positions are consecutive, put the three
positive positions at $0,1,2$. The base count is $f(D)-(D-1)$.
No outside mark sees the omitted cycle edge. A mark from $z$ to
$U_1'$ would also miss that edge, so $h_z\le D-1$ and $z$ marks
to $U_0'$ or $U_2'$. A mark to $U_0'$ forces $z$ to see every
vertex of $U_2'$ having a marked neighbor in $U_3'$ by comparison
with the marked edges to $U_3'$, where $s'_3=0$. If such a
neighbor is absent at some vertex of $U_2'$ then $L'\ge d+1$
and \eqref{eq:multi-extra-capacity} suffices. Otherwise the
adjacency to $U_2'$ is complete. If these pairs are all blue the
cut costs at least $d$. If one is marked, reflection forces
complete adjacency to $U_0'$ too. The two classes now have no
remaining extra capacity and the middle class cannot have a
same-class mark. The only extra marks meet $z$ and number at most
$D-1$. In every subcase $\mu\ge d$ suffices, since its fractional
part is $\frac{1}{4}$.

Suppose now that there is no vertex with three or more positions.
Then there are no vertices of signature size zero. If all outside vertices have
two positions, the budget and parity give $R-K=1$ and there are
$n=\frac{R-1}{2}$ outside vertices. The total number of unmarked host
edges is at most $\frac{D-5}{4}$. Indeed, summing the differences
between $D$ and the marked degree over the $\frac{5D-1}{2}$ vertices
counts at least twice the number of unmarked host edges.
Hence $\sum b_x\le \frac{D-5}{4}-L<a_*-\sqrt L$,
using $a_*\ge\frac{D-1}{4}$. Every outside vertex has a marked core
neighbor by the vertex-cover bound, while the comparison for
consecutive positions excludes those positions. Apply the enlargement
rule to every eligible vertex. Two remaining vertices cannot be joined
by a marked edge.
Their combined old-core neighborhood would have to omit at most
$\sqrt L$ points in at least three classes, since two consecutive
omitted sets have product at most $L$. In any class outside the
intersection of their two designated class sets, their combined
neighborhood omits at least $a_*-b_x-b_y>\sqrt L$ points. The
intersection has at most two classes. If no vertex remains, $R'=1$.
Otherwise take a marked edge to $U_i'$ from a remaining vertex $x$
of positions $\{i,j\}$. Each original vertex of signature size two pays
one blue anchor pair to $L'+L_{\rm cut}$. For $x$ choose its pair
to $c_i$, which is disjoint from its $b_{x,j}$ pairs. By
\eqref{eq:multi-rectangle} and \eqref{eq:multi-extra-capacity},
$\mu\ge n+a_*=\frac{D-1}{2}$, as required.

It remains to consider $a>0$ vertices of signature size one and no
vertex of signature size three. If $a$ is even then $R-K=1$ and the loss
is at least $\frac{D}{2}+R-1+a(t+1)>D-\frac{5}{2}$. Hence $a$ is odd and $R=K$.
Apply the enlargement rule and let $n$ be the number of remaining
vertices of signature size two, so $R'=a+2n$. After subtracting the base
losses, \eqref{eq:multi-loss-budget} implies
\[
 2L+\sum_{\rm remaining\ two}(\ell_x-2)
 \le T=2a_*-\frac{5}{2}-a(t+1).
\]
A vertex with $b_x\ge a_*$ contributes at least $2a_*-2>T$ to the
left side. For a vertex with $b_x<a_*$ and consecutive signature
positions, or a vertex without a marked edge to the core, its
contribution together with $2L$ is at least $2a_*$ by the bounds
$a_*+1-\sqrt L$ and $q+1$, respectively. Therefore each remaining
vertex with $k_x=2$ has $b_x<a_*$, nonconsecutive signature positions,
and a marked edge to the core. Its two anchor pairs are blue.

If two such vertices $x,y$ mark to one another, their signatures
are distinct. Equation \eqref{eq:multi-low-blue-outside} shows that
their marked-neighbor counts in the old core sum to at most
$2R'-n$, since any shared position has $s_i'\le a+n$.
Consequently
\[
 b_x+b_y\ge D+q-2R'+n\ge a_*+2t+n-\frac{L}{a_*}.
\]
The budget gives $b_x+b_y\le \frac{T}{2}-n+4-L$. These contradict for
$n\ge2$, $a\ge1$ and $t\ge-\frac{1}{2}$.

Suppose $xy$ is marked, where $k_x=1$ and $y$ is a remaining
vertex with $k_y=2$. After a rotation of the cycle, the signature
of $x$ is $\{0\}$ and that of $y$ is $\{1,3\}$ or $\{2,4\}$. The omitted anchor stars require
complete combined adjacency to the three indicated old classes.
If $x$ has a marked edge to $B$, it has at least $a_*+1$ blue pairs
in $U_0$. The designated class $U_1$ or $U_4$ for $y$ is short of
at least $a_*-b_y$ of its vertices. These must be covered by
additional blue pairs of $x$, since $x$ only marks to $U_2$ or
$U_3$ in the old core. If $x$ has no marked edge to $B$, the
three required classes have at least $D+a_*$ vertices and $y$
has at most $R'$ marked neighbors there. Both cases imply
$b_x+b_y\ge2a_*+1$. The budget then requires at least
$D+a(t+1)-2>D-\frac{5}{2}$ in losses after including the other vertices.
There are also no marks between two vertices of signature size one.
Thus $f=0$. If $n>0$, then \eqref{eq:multi-rectangle} and the blue edges
from one vertex of signature size one to $B$ give $I+L_{\rm cut}\ge2a_*+1>\frac{D}{2}$.
If $n=0$ and $a\ge3$ then $L_{\rm cut}\ge3(a_*+1)>\frac{D}{2}$.
The last possibility is $R'=a=1$, already settled.


\subsubsection{Small-degree slack thresholds and a single repeated pair}
\label{app:multi-small-slack}

For $(D,R)=(11,7),(15,9)$, write $D=4k+3$ and $R=2k+3$.
Here $k=2,3$ and $a_*=k$. Equation \eqref{eq:multi-rough-discharge}
gives $L\le9,13$, respectively. For $L<D$, the old marked
vertex-cover number satisfies the improved bound
\begin{equation}\label{eq:multi-small-cover}
 \tau\ge D+a_*-\max\{0,\lfloor\sqrt L\rfloor-1\}.
\end{equation}
To prove it, consider an independent set of size $\alpha>D$ in
the marked old core, with $x_i$ points in $U_i$. It occupies a
consecutive pair of positions, since any two nonconsecutive classes
have total size at most $D$. If the occupied positions induce a
connected graph, summing $x_i x_j\ge x_i+x_j-1$ on a spanning tree
gives $L\ge\alpha-1\ge D$. If they induce a disconnected graph,
the occupied positions consist of one edge and one isolated point.
The two classes of the edge cannot contribute their own anchors,
because the anchor spokes to the adjacent class are marked. Either
of these classes and the isolated class now contribute at most
$D-1$ points. Hence the adjacent classes
each contribute at least $\alpha-D+1$ and $L\ge(\alpha-D+1)^2$.
Subtracting the upper bound for $\alpha$ from $|B|$ proves
\eqref{eq:multi-small-cover}.

Use $q=\tau_{\rm lo}-D$ with $\tau_{\rm lo}$ the right-hand side
of \eqref{eq:multi-small-cover} in the sharper discharging.
Initially $q\ge0$. A vertex with $k_x=1$ has loss at least
$\frac{1}{2}$, and one with $k_x=2$ has loss at least $2$. All other
losses are nonnegative. Thus $L\le 2k$ and $R-K\in\{0,1\}$. For $c\ge2$
potential recipients, positive transfer requires
$h_x\le c(Q_x-q)$ and $Q_x>\frac{cq}{c-1}$. Substituting the integer
values $Q_x<\frac{D}{2}$ bounds the number of recipients as follows.

For $D=11$, two recipients with $L\le3$ have $q=2$. The only
possible $Q_x$ for a positive transfer is $5$, so the rate
is at most $\frac{1}{6}$ and each recipient can receive on at most
three edges. Their combined loss is at least $10>\frac{17}{2}$.
At $L=4$ even one recipient's loss with $2L$ is at least
$\frac{11}{2}-\frac{81}{20}+8>\frac{17}{2}$. For $D=15$, three recipients with
$L\le3$ have $q=3$ and each can receive at most $\frac{18}{7}$.
Their combined loss is greater than $\frac{25}{2}$. When $L\ge4$,
$q=2$ and the recipient bound is $\frac{33}{8}$ each, again exceeding
$\frac{25}{2}$ with the $2L$ term. With two recipients and $q=3$,
only $Q_x=7$ can give a positive transfer, at a rate of at most
$\frac{1}{8}$ for at most four edges per recipient. With two recipients
and $q=2$, each can receive at most $\frac{7}{2}$, and $L\ge4$.
Both violate the budget, so there is at most one vertex of signature size
three in either degree, with no positive transfer.

We give the remaining cases for $D=11,15$ in terms of
$d=\frac{D-1}{2}=2k+1$. A contradiction follows from $\mu\ge d$.
We use two consequences of enlargement. First, each absorbed vertex
$x$ contributes at most $D-h_x$ new blue consecutive pairs and at
least one such pair, since one of its two anchor pairs is blue.
Thus, if $p$ vertices are absorbed,
\begin{equation}\label{eq:multi-new-blue}
 L+p\le L'\le L+\sum_{x\ {\rm absorbed}}(D-h_x).
\end{equation}
The vertex cover number of the enlarged marked core is at least
$\tau'=D+a_*'-\frac{L'}{a_*'}$, rounded up to the next integer,
where $a_*'=\frac{D-R'}{2}$. Second, if all $f$ outside marked
edges meet one vertex and that vertex has at most $p$
distinct outside marked neighbors, then
\begin{equation}\label{eq:multi-center-cut}
 f\le\frac{p}{p+1}(2D-\tau').
\end{equation}
Indeed, one of its neighbors has outside marked degree at least
$\frac{f}{p}$, and the outside marked degrees of the two ends sum to
at most $2D-\tau'$. The bound is immediate when $p=0$.

Suppose there is one vertex $z$ of signature size three. Write $a$
for the number of vertices of signature size one. If $R-K=1$, the
unused anchor capacity and the loss at $z$ already sum to $D$, so
$R=K$ and $a$ is even. The $k-\frac{a}{2}$ vertices of signature size
two, the $a$ vertices of signature size one, and $z$ give total
loss at least $4k+\frac{3}{2}-\frac{a}{2}$. After these losses,
\eqref{eq:multi-loss-budget} allows at most $\frac{a}{2}-1$ more.
It follows that $a=2,4$, or that $a=6$ and $k=3$.

For $a=2,4$ we have $L=0$ and every vertex of signature size two
has $b_x=1$. These vertices have nonconsecutive positions and
enter the enlargement. Indeed, for consecutive positions a marked
edge into the enlarged core requires $b_x\ge k+1$, and an outside
marked edge with no enlarged core mark requires $b_x\ge q+1=k+1$.
A vertex of the nonconsecutive type left outside would have two
blue anchor pairs. When $a=6$ there are no vertices of signature
size two. All outside marks after enlargement meet $z$. Every
vertex of signature size one has at least $k+1$ blue pairs into
$B$, since a vertex without the two listed kinds of marked neighbor
would lose $\frac{D}{2}$, exceeding the available excess loss
$\frac{a}{2}-1$. If a vertex of signature size zero marks to $z$,
its entire core neighborhood is blue and has size at least
$\tau'-D+f$. In these cases $\tau'\ge D$, by
\eqref{eq:multi-new-blue} and the bounds on $L'$ below. The
capacity inequality \eqref{eq:multi-extra-capacity} thus gives
$\mu\ge a(k+1)\ge d$. We may assume that all outside marked
edges join $z$ to vertices of signature size one.

If $a=2$, both vertices of signature size one attain their
minimum loss. Each has marked degree $3k+2$ and exactly $k+1$
unmarked incidences, all blue pairs into $B$. The absorbed
vertices have $D-h_x=1$, so $R'=5$, $a_*'=2k-1$, and
$L'\le k-1$. Each vertex of signature size one has an old core
mark. Otherwise its old core degree would be $k+1$, and every
outside marked neighbor would have at least $D-1$ old core
neighbors by the vertex cover bound with $q=k$. Each could give
only one outside marked copy, and there are at most $k$
possible marked neighbors of this kind.

For a vertex of signature $\{j\}$ and a marked neighbor in
$U_i'$, where $i=j\pm2$, use the comparison with the anchor
star from the opposite side of $c_j$. At least
$a_j'-s_i'+1-L_{\star}\ge2k-L_{\star}$ vertices of $U_j'$
are blue neighbors of the vertex of signature size one. Here
$L_{\star}$ counts the blue pairs of that star into $U_j'$.
As it has only $k+1$ blue neighbors, $L_{\star}\ge k-1$.
Stars for different choices of $j$ or different choices of
$i$ share at most a red edge on the anchor cycle, so their
blue pairs are disjoint. It follows from $L'\le k-1$ that the
two vertices have the same signature and mark to the same
class only. The index of this class is used at most once by
the remaining signatures, namely by $z$. Since $R'=K'$,
$s_i'\le1$. Their total marked degree is consequently at most
$a_i'+f$. By the rounded cover bound,
$\tau'\ge D+2k-1$, so \eqref{eq:multi-center-cut} gives
$f\le\frac{2(2k+4)}{3}$. As $a_i'\le2k+4$, this sum is
strictly less than $6k+4$, a contradiction.

If $a=4$, then $R'=7$, and \eqref{eq:multi-new-blue} gives
$L'\le k-1$. Hence $\tau'\ge D+2k-2$ and
$f\le\frac{4(2k+5)}{5}$, which gives $f\le7$ for $k=2$ and
$f\le8$ for $k=3$. In both cases $L_{\rm cut}\ge4(k+1)$
and \eqref{eq:multi-extra-capacity} gives $\mu\ge d$.
If $a=6$, then $k=3$, $L_{\rm cut}\ge24$ and $f\le15$,
which again suffice.

Suppose there is no vertex of signature size at least three.
There are no vertices of signature size zero, and every vertex
of signature size one has at least $k+1$ blue pairs into the
old core. If $a$ is even, then $R-K=1$, and the base losses
require $a\ge6$, leaving at most $\frac{a}{2}-3\le1$ more.
Enlarge the core by all vertices with nonconsecutive signatures
and $b_x<k$ that are complete to both designated classes.
Each remaining vertex of signature size two with $b_x\ge k$
or with two blue anchor pairs costs at least two beyond the
base loss of two. A remaining vertex with consecutive positions
and a core mark, or without any core mark, costs at least $2k$
beyond its base loss when $2L$ is included, by the bounds
$k+1-\sqrt L$ and $q+1$ for $b_x$. Thus all vertices of
signature size two are absorbed. There are no outside marks and
$L_{\rm cut}\ge a(k+1)$.

If $a$ is odd, then $R=K$. Let $n$ count the remaining vertices
of signature size two after enlargement, so $R'=a+2n$. After
subtracting the base losses, the remaining budget is
\begin{equation}\label{eq:multi-small-odd-budget}
 2L+\sum_{\substack{x\ \text{remaining}\\ k_x=2}}(\ell_x-2)
 \le T=2k-\frac{5}{2}+\frac{a}{2} .
\end{equation}
In the following count, a remaining vertex of signature size two
that has nonconsecutive positions, $b_x<k$, and a marked edge
into the enlarged core has two blue anchor pairs. By
\eqref{eq:multi-low-blue-outside}, its old core degree is at
most $R'+b_x$. A marked edge between such a vertex $y$ and a
vertex $x$ of signature size one requires
$b_x+b_y\ge2a_*+1=2k+1$, by the comparison with the two
omitted anchors in \textbf{Appendix~\ref{app:multi-half}}.

If $a\ge5$ and $n\le1$, then all outside marks meet the sole
remaining vertex of signature size two, if it exists. For $n=0$,
$f=0$. For $n=1,k=2$, $L_{\rm cut}\ge5(k+1)+1=16$.
For $n=1,k=3,a=5$ we have $L\le3$, so a remaining vertex
of signature size two has at least two blue cut pairs, whether
its signature is consecutive or nonconsecutive or it has no
core mark. This gives $L_{\rm cut}\ge22$. If $a\ge7$ and
$k=3$, the vertices of signature size one alone give at least
$28$. In each case $f\le D$ and the capacity inequality suffices.

If $a\ge5$ and $n\ge2$, then $k=3$, $a=5$, $n=2$, and
$T=6$. The two remaining vertices of signature size two have
$b_x,b_y\ge2$. If either has $b_x<3$ or $b_y<3$, it has
nonconsecutive positions and a mark into the enlarged core by
the lower bounds on $b_x,b_y$ for the other types. If the two
vertices have a marked edge between them, the sum of their
outside marked degrees is at most $D-q$ by the old core cover
bound, so $f\le D$. Otherwise, if both have blue degree two,
they use at least four units of the budget $T$. At most one
vertex of signature size one can use the remaining two units
to increase its blue degree from four to five. The comparison
$b_x+b_y\ge7$ forces all outside marked copies to meet that
vertex, so $f\le D$. If instead one of the two vertices has
blue degree at least three, the two vertices exhaust $T$.
The vertices of signature size one all have blue degree four,
preventing a marked edge to the remaining vertex of blue
degree two. Again $f\le D$. In each case the blue cut has at
least 24 edges.

If $a=3$, then $T=2k-1$. Each remaining vertex of signature
size two has blue degree at least two, by the same comparisons.
For $n=2$ we must have $k=3$, $L=0$, and both of those vertices
have blue degree two and nonconsecutive positions with a marked
edge to the enlarged core. All three vertices of signature size
one have blue degree four. None can mark to either remaining
vertex of signature size two. If the two vertices of size two
have a marked edge, their signatures are distinct and share at
most one position, of slack at most $a+n$. By
\eqref{eq:multi-low-blue-outside} their combined old core degree
is at most $2R'-n+4=16<D+q$, also impossible. Hence $f=0$.
For $n=1$ we have $R'=5$, and
\eqref{eq:multi-new-blue} and \eqref{eq:multi-small-odd-budget}
give $L'\le(k-1)+T-2-L\le3k-4$. The resulting cover bounds
are $\tau'\ge14$ for $k=2$ and $\tau'\ge19$ for $k=3$.
Equation \eqref{eq:multi-center-cut} gives $f\le6$ and $f\le8$,
respectively. The blue cut has at least 11 or 14 edges,
respectively, and in the latter case $L'\ge2$ by
\eqref{eq:multi-new-blue}. For $n=0$ there are no outside marks,
so the blue pairs from the three vertices of signature size
one suffice.

Finally, let $a=1$, so $T=2k-2$. The lower bounds for vertices
with consecutive positions or with no enlarged core mark
exclude them when their blue degree is less than $k$. If a
remaining vertex $y$ has $b_y\ge k$, its loss above two uses
all of $T$. It is the only remaining vertex of signature size
two, with $b_y=k$ and $h_y=D-k$. We have $L=0$, there are $k$
absorbed vertices, and $L'\le k$. The vertex $x$ of signature size
one has $b_x=k+1$ and $h_x=D-b_x$. Now $R'=3$. If $f=0$, the
blue cut has at least $2k+1=d$ edges. If $f>0$, all outside
marks run between $x$ and $y$. Rotate so that their signature union
is $\{0,2,3\}$. The three corresponding slacks are one and all
other slacks are zero. Put $A_i=U_i'$ for $i=0,2,3$. The sizes
are $d+1,d,d$, respectively. Let $w$ count their vertices
with no red pair to either omitted anchor. These missing red
pairs give distinct blue base pairs, so $w\le L'$. Write
$Q_x',Q_y'$ for the enlarged core degrees. The omitted anchors
force the union of the two neighborhoods to cover at least
$3d+1-w$ vertices of $A_0\cup A_2\cup A_3$. Each end has at most
$D-f$ core neighbors, so $Q_x',Q_y'\ge d+f-w$ and
$2f\le d+1+w$. If one end has no core mark, its blue cut pairs
and $L'\ge w$ give $\mu\ge d$.

Assume both ends have a core mark. If the signature of $x$
is $\{0\}$, then $x$ can only mark into $A_2$ or $A_3$.
Take a marked neighbor in $A_2$, reflecting if necessary.
Its unique extra degree unit is used by this edge, so the
red star from $c_4$ into $A_0$ forces $x$ to see each of its
red neighbors in $A_0$, all through blue pairs. If the star
has $w_0$ blue pairs, this gives at least $d+1-w_0$ blue cut
pairs at $x$. The only available core marks at $y$ go to
$A_0$, by comparison with the cycle edges. At most $w_0$
vertices there still have a degree unit free for $y$. Thus
$L'+L_{\rm cut}-f\ge L'+d+1-w_0+Q_y'-w_0-f\ge2d+1-2w_0\ge d$,
using $w_0\le L'\le k$ and $d=2k+1$.

Otherwise, reflect to give $x$ signature $\{2\}$. It can
mark only into $A_0$. Choose a marked neighbor $u\in A_0$.
The red star from $c_3$ into $A_2$ forces at least $d-w_2$
blue cut pairs from $x$ into $A_2$, where $w_2$ is the
number of blue pairs of that star. If $uc_1$ is red, comparing
that edge with any proposed mark from $y$ into $A_3$ excludes
all such marks. Indeed, neither $y$ nor the vertex in $A_3$
sees $c_1$, and $u$ has used its only extra unit on $xu$.
Therefore all pairs from $x,y$ into $A_3$ are blue. They
cover at least $d-w_3$ vertices there, where $w_3$ counts the
vertices in $A_3$ without a red pair to an omitted anchor.
The corresponding blue base pairs in $A_3$ are disjoint
from the $w_2$ pairs of the star, so $L'\ge w_2+w_3$.
This gives $\mu\ge 2d-f\ge d$, as $2f\le d+1+L'\le2d$.
If every choice of $u$ instead gives a blue pair $uc_1$,
each marked copy from $x$ into $A_0$ gives a distinct
blue base pair. Now $L'\ge h_x-f$, and the old blue pairs
at $x,y$ give $\mu\ge h_x+b_x+b_y-2f=D+k-2f\ge d$.

If all remaining vertices have $b_x<k$, then $n>0$ is possible
only for $k=3$, with $n\le2$ and each blue degree equal to two.
For $n=1$ an outside mark would force $R'=3$ with three
positive slacks all equal to one. The bound
$b_x\ge a_*+1-s_i'$ preceding \eqref{eq:multi-rectangle}
would then give blue degree at least three. For $n=2$, the
two vertices use all of $T$, so $L=0$ and the vertex of
signature size one has blue degree four. Neither vertex can
mark to it, and their combined old core degree is at most
$2R'-n+4=12<D+q$ if they mark to one another. Thus $f=0$.
If $n>0$, apply \eqref{eq:multi-rectangle} to one remaining
vertex and add the $a_*+1$ blue cut pairs of the vertex of
signature size one. These incidences are disjoint and give
$\mu\ge2a_*+1=d$. If $n=0$, then $R'=1$ and only the base
count is needed. This completes the two cases $D=11,15$.

For $(D,R)=(13,9)$, \eqref{eq:multi-joint} gives
\[
 L\le \mu+\frac{1}{2}\,\frac{1+r}{2+r}R(R-\frac{D}{2})<13,\qquad r=\frac{9}{13} .
\]
It follows from \eqref{eq:multi-small-cover} that $q\ge0$.
For a vertex with $k_x=1$, the initial lower bound on its loss is $-\frac{1}{2}$, so
\eqref{eq:multi-loss-budget} gives $R-K\le2$ and $L\le7$ and
hence $q\ge1$. For completeness, we record the transfer bounds
that reduce this case to the terminal counts in the proof supplement.
If there are $c$ recipients, then $c\le3$, and a sender with positive
contribution has $Q<h\le\min\{13-Q,c(Q-q)\}$. It follows that
$Q\le6$ and that a recipient has outside marked degree at most
$6-q$. With $c=3$, there are no vertices of signature size one.
For $q=1$, each marked edge transfers at most $\frac{5}{9}$, so the
three recipients lose at least $3(\frac{13}{2}-\frac{25}{9})>\frac{21}{2}$.
For $q=2$, each marked edge transfers at most $\frac{3}{8}$, so the
loss is at least $15$. For $c=2,q=1$, each marked edge transfers at
most $\frac{3}{8}$. Here $L\ge4$ and there are at most three vertices
of signature size one, giving combined loss and blue deductions
at least $8+13-2\cdot\frac{15}{8}-\frac{3}{2}>\frac{21}{2}$.
These rates follow by checking $Q\in\{1,\ldots,6\}$ in the minimum
above. If $c=2,q=2$, the only possible pairs $(Q,h)$ are
$(5,6)$ and $(6,7)$, and a positive sender meets both recipients.
They have at most eight outside marked incidences in total, so
there is at most one positive sender, transferring one unit.
The loss bound then forces $L=R-K=0$, three vertices of signature
size one with $h_x=10,b_x=3$, and two recipients of signature
size three with $h_x=13,b_x=0$. The sender leaves at most one
outside marked incidence at the recipients for the three vertices
of signature size one and other neighbors. The three vertices
therefore contribute at least 29 marked copies across the old
core cut, and the recipients contribute at least 18. The sender
and the three vertices of signature size one also contribute at
least $5+9$ blue pairs across that cut. The total is at least 61,
whereas \eqref{eq:multi-core-identities} bounds the capacity of
the cut by $18+2\lfloor\frac{81}{4}\rfloor=58$.

There is thus at most one recipient, with no positive transfer.
The exact integer comparisons and remaining cases are proved in
Sections 2--5 of \emph{The exact bound for degree thirteen with
slack nine} in the proof supplement indexed by \path{PROOF_NOTES.md}.
Section 2 supplies
the cover bounds for the enlarged core, and Sections 3--5 split
according to the number of vertices of signature size one and
the signatures remaining after enlargement. Each case gives
$\mu\ge6$ or the base count for $R'=1$. Consequently $R\ge11$
for degree thirteen. For degree nine, the finite reductions
in \textbf{Appendix~\ref{app:multi-nine}} handle $R=7,9$.

We also need to exclude $b=1$ after taking exactly $m=f(D)+1$
marked copies. The projection to distinct host pairs preserves
both the strong-clique obligation and the anchor classes. If
$b=1$, this projection has exactly $f(D)$ marked pairs. In the
simple equality configuration with class sizes
$(d,d,d+1,d+1,d)$, where $d=\frac{D-1}{2}$, every marked pair has
an endpoint of host degree at least $D$ in the projection.
A second marked copy on it contradicts the host capacity.

The simple equality test in the companion files uses a cycle with
largest old core and identifies the five slack entries under
dihedral symmetry. For $D=9,11$, the only possible exception
to the simple equality configuration is $(1,1,2,1,2)$.
For $D=13,15$, the exceptions can only be $(1,1,2,1,2)$,
$(1,1,3,1,3)$ or $(1,2,2,2,2)$. The last two have total
slack nine. Thus the preceding half-degree and special
slack bounds rule out every exceptional case at $D=11,13,15$
in the multigraph counterexample. At $D=9$, the case $b=1$ with a
triangle-free marked graph in the simple support is ruled out by the fixed
$27$-formula replay in \path{d9_anchor_cases.py}.
With a marked triangle the order is $24$ or $25$ by the
degree-nine order reduction. The only exceptional shape is
$(1,1,2,1,2)$, and the replay with \path{d9_pentagon.py},
\path{--bundles 1 --triangles require}, and
\path{--select s11212} checks its fixed unsatisfiability
proof. This is the case \path{d9_b1_triangle/s11212.json}
in the companion archive. The precise reduction and replay
transcripts are indexed under \emph{A single repeated marked pair}.

